\chapter{Euclidean Distances Between Function Values}

This chapter investigates the geometric distances between consecutive
integer points on the curve $k(n) = \frac{1}{6}(2n^2 + 3n + 1)$.

\section{Motivation}

The function takes integer values for all $n \equiv 1, 5 \pmod{6}$. This
property allows us to consider discrete points $(n, k(n))$ on the curve.
A natural question is: \textbf{What are the Euclidean distances between
consecutive integer points?}

\section{Point Distances in the Plane}

\subsection{Definition}

For two integer arguments $n_1, n_2$ with $n_1, n_2 \equiv 1, 5 \pmod{6}$,
the points $(n_1, k(n_1))$ and $(n_2, k(n_2))$ lie on the curve.

\begin{definition}[Euclidean Distance Between Curve Points]
The Euclidean distance between two points on the curve is:
\begin{equation}
d(n_1, n_2) = \sqrt{(n_2 - n_1)^2 + (k(n_2) - k(n_1))^2}
\end{equation}
\end{definition}

\subsection{Numerical Examples}

\begin{example}[Distances between integer points]
The following table shows the Euclidean distances for the first
integer-valued pairs:

\begin{table}[H]
	\centering
	\caption{Euclidean distances between consecutive points of
		the sequence $n \equiv 1,5 \pmod{6}$}
	\begin{tabular}{cccccc}
		\toprule
		$(n_1, n_2)$ & $\Delta n$ & $k(n_1)$ & $k(n_2)$ & Distance $d$ & Exact value \\
		\midrule
		$( 1,  5)$ & 4 &   1 &  11 &  10.7703296143 & $\sqrt{116}$  \\
		$( 5,  7)$ & 2 &  11 &  20 &   9.21954445729 & $\sqrt{85}$   \\
		$( 7, 11)$ & 4 &  20 &  46 &  26.3058928759 & $\sqrt{692}$  \\
		$(11, 13)$ & 2 &  46 &  63 &  17.1172427686 & $\sqrt{293}$  \\
		$(13, 17)$ & 4 &  63 & 105 &  42.1900462195 & $\sqrt{1780}$ \\
		$(17, 19)$ & 2 & 105 & 130 &  25.0798724080 & $\sqrt{629}$  \\
		$(19, 23)$ & 4 & 130 & 188 &  58.1377674150 & $\sqrt{3380}$ \\
		$(23, 25)$ & 2 & 188 & 221 &  33.0605505096 & $\sqrt{1093}$ \\  % 25 = 5² composite
		$(25, 29)$ & 4 & 221 & 295 &  74.1080292546 & $\sqrt{5492}$ \\
		$(29, 31)$ & 2 & 295 & 336 &  41.0487515035 & $\sqrt{1685}$ \\
		\bottomrule
	\end{tabular}
\end{table}
\end{example}

\section{Closed-Form Formulae}

\subsection{Exact Distance Formula}

\begin{proposition}[Closed-form formula for the distance]
With $k(n) = \frac{1}{6}(2n^2 + 3n + 1)$ and $\Delta n = n_2 - n_1$:
\begin{align}
k(n_2) - k(n_1) &= \frac{\Delta n}{6}(4n_1 + 2\Delta n + 3)
\end{align}

Hence:
\begin{equation}
d(n_1, n_1 + \Delta n) = \Delta n \cdot \sqrt{1 + \frac{1}{36}(4n_1 + 2\Delta n + 3)^2}
\end{equation}
\end{proposition}

\begin{proof}
Direct computation:
\begin{align}
k(n_2) - k(n_1) &= \frac{1}{6}[2(n_2^2 - n_1^2) + 3(n_2 - n_1)]\\
&= \frac{1}{6}[2(n_2 - n_1)(n_2 + n_1) + 3(n_2 - n_1)]\\
&= \frac{\Delta n}{6}[2(2n_1 + \Delta n) + 3]\\
&= \frac{\Delta n}{6}(4n_1 + 2\Delta n + 3)
\end{align}

Substitution into the distance formula yields the claim.
\end{proof}


\subsection{Special Cases}

% ── Δn = 2 and Δn = 4 side by side ─────────────────────────
\noindent
\begin{minipage}[t]{0.48\textwidth}

\subsubsection{Case $\Delta n = 2$}

\begin{equation}
d(n, n+2) = 2\sqrt{1 + \frac{1}{36}(4n + 7)^2}
\end{equation}

\begin{example}[Numerical values for $\Delta n = 2$]
\begin{table}[H]
\centering
\begin{tabular}{cc}
\toprule
$n_1$ & Distance $d$ \\
\midrule
5  & 9.2195444573  \\
11 & 17.1172427686 \\
17 & 25.0798724080 \\
23 & 33.0605505096 \\
\bottomrule
\end{tabular}
\end{table}
\end{example}

\end{minipage}
\hfill
\begin{minipage}[t]{0.48\textwidth}

\subsubsection{Case $\Delta n = 4$}

\begin{equation}
d(n, n+4) = 4\sqrt{1 + \frac{1}{36}(4n + 11)^2}
\end{equation}

\begin{example}[Numerical values for $\Delta n = 4$]
\begin{table}[H]
\centering
\begin{tabular}{cc}
\toprule
$n_1$ & Distance $d$ \\
\midrule
1  & 10.7703296143 \\
7  & 26.3058928759 \\
13 & 42.1900462195 \\
19 & 58.1377674150 \\
\bottomrule
\end{tabular}
\end{table}
\end{example}

\end{minipage}
% ─────────────────────────────────────────────────────────────

\subsubsection{Case $\Delta n = 6$}

\begin{equation}
d(n, n+6) = 6\sqrt{1 + \frac{1}{36}(4n + 15)^2}
\end{equation}

\begin{remark}
The case $\Delta n = 6$ corresponds to a full period in the congruence
structure modulo~$6$. Geometrically, this means a jump across both
congruence classes.
\end{remark}

\section{Comparison: Distance vs.\ Arc Length}

The Euclidean distance $d$ between two points is always less than or equal to
the arc length $L$ along the curve, since the straight line represents the
shortest possible path. This is a fundamental metric property.

\begin{proposition}[Distance as a lower bound]
For $n_1 < n_2$ it always holds that:
\begin{equation}
d(n_1, n_2) \leq L[n_1, n_2]
\end{equation}
where $L[n_1, n_2] = \int_{n_1}^{n_2} \sqrt{1 + [k'(x)]^2}\, dx$ is the arc length.
\end{proposition}

\begin{example}[Comparison: distance vs.\ arc length]
\begin{table}[H]
\centering
\caption{Euclidean distance compared with arc length}
\begin{tabular}{ccccc}
\toprule
Interval & Distance $d$ & Arc length $L$ & Ratio $d/L$ & $L - d$ \\
\midrule
$[1, 5]$   & 10.7703296143 & 10.8394091489 & 0.9936270018 & $6.908 \cdot 10^{-2}$ \\
$[5, 7]$   &  9.2195444573 &  9.2210749306 & 0.9998340244 & $1.530 \cdot 10^{-3}$ \\
$[7, 11]$  & 26.3058928759 & 26.3101623759 & 0.9998377243 & $4.270 \cdot 10^{-3}$ \\
$[11, 13]$ & 17.1172427686 & 17.1174799293 & 0.9999861451 & $2.372 \cdot 10^{-4}$ \\
$[13, 17]$ & 42.1900462195 & 42.1910659377 & 0.9999758309 & $1.020 \cdot 10^{-3}$ \\
\bottomrule
\end{tabular}
\end{table}

\textbf{Observation:} The ratio $d/L$ is in all cases $\leq 1$, as
guaranteed by the proposition. For short intervals ($\Delta n = 2$), distance
and arc length are nearly identical, since the curve is locally almost
rectilinear. The deviation is somewhat larger for small $n$ (more strongly
curved region) and approaches the value~$1$ for large $n$.
\end{example}

\section{Distance Sequence}

Consider the sequence of all consecutive distances:

\begin{definition}[Distance Sequence]
Let $\{n_k\}$ be the ascending sequence of all integer arguments of $k(n)$
(i.e.\ $n_k \equiv 1, 5 \pmod{6}$, cf.\ Chapter~\ref{sec:arithmetik}).
Then we define:
\begin{equation}
\{d_k\} = \{d(n_k, n_{k+1})\}_{k=1}^{\infty}
\end{equation}
\end{definition}

\begin{example}[The first terms]
With $n_1 = 1, n_2 = 5, n_3 = 7, n_4 = 11, \ldots$ we obtain:
\[
\{d_k\} = \{10.7703,\ 9.2195,\ 26.3059,\ 17.1172,\ 42.1900,\ 25.0799,\ \ldots\}
\]
\end{example}

\begin{remark}[Irregularity]
The sequence $\{d_k\}$ is not monotone, since the spacings
$\Delta n_k = n_{k+1} - n_k$ alternate between $2$ and $4$. However, the
sequence exhibits a clear growth pattern: both the local maxima and the local
minima increase monotonically with $k$.
\end{remark}

\section{Summary}

\begin{table}[H]
\centering
\caption{Distance properties by $\Delta n$}
\begin{tabular}{ccc}
\toprule
$\Delta n$ & Type & Distance formula \\
\midrule
2 & Adjacent values (minor) & $2\sqrt{1 + \frac{(4n+7)^2}{36}}$ \\
4 & Congruence class change (major) & $4\sqrt{1 + \frac{(4n+11)^2}{36}}$ \\
6 & Full period & $6\sqrt{1 + \frac{(4n+15)^2}{36}}$ \\
\bottomrule
\end{tabular}
\end{table}

\begin{remark}[Significance]
The distance analysis shows that the geometric spacings between integer
function values exhibit a clear pattern directly linked to the congruence
structure modulo~$6$. For large $n$, all distances grow proportionally to $n$,
with the proportionality factor depending on $\Delta n$.
\end{remark}
